\section{MAT-freeness implies strong chordality}
\subsection{MAT-labeling of graphs}
\label{subsec:MAT-labeling of graphs}
 
In this subsection, we show that MAT-free graphic arrangements (Definition \ref{def:MAT-partition-free}) can be completely determined by a  special edge-labeling of graphs. 
First we show that the condition of being admitted a ``partition'' of a (nonempty) MAT-free arrangement is actually implied by the three conditions \ref{definition MAT partition 1}, \ref{definition MAT partition 2}, and \ref{definition MAT partition 3}. 
\begin{proposition}\label{no internal empty block}
An arrangement $ \mathcal{A} $ is MAT-free if and only if $ \mathcal{A} $ can be decomposed into a disjoint union of \emph{possibly-empty} subsets $ \pi_{1}, \dots, \pi_{n} $ of $ \mathcal{A} $ satisfying \ref{definition MAT partition 1}, \ref{definition MAT partition 2}, and \ref{definition MAT partition 3}. 
\end{proposition}
\begin{proof}
If $ \A = \varnothing $, then the statement is clear. Suppose $ \A \ne \varnothing $. 
We only need to show the ``if'' part, namely, the existence of a disjoint union of  possibly-empty subsets~$ \pi_{1}, \dots, \pi_{n} $ of $ \mathcal{A} $ satisfying \ref{definition MAT partition 1}, \ref{definition MAT partition 2}, and \ref{definition MAT partition 3} implies the existence of an MAT-partition of $\A$.

First we show that $ \pi_{1} \neq \varnothing $. 
Suppose to the contrary that $ \pi_{1} = \varnothing $. 
Thus $n\ge2$.
Let~$ k \geq 2 $ be the minimal integer such that $ \pi_{k} \neq \varnothing $. 
Then \ref{definition MAT partition 3} yields $ 0 = |\mathcal{A}_{k-1}| \geq k-1 \geq 1 $, a contradiction. 
Thus $ \pi_{1} \neq \varnothing $. 

Let $p:= \max \{ 1 \le i\le n \mid \pi_i \ne\varnothing \}$.
We will show that  $\pi_{k} \neq \varnothing $ for all $ k \in [p] $ which in turn implies that $ (\pi_{1}, \dots, \pi_{p}) $ is an MAT-partition of $ \mathcal{A} $. 
Suppose to the contrary that there exists   $2\le  k< p$ such that $ \pi_{k} = \varnothing $ and choose minimal such $k$.
Set $ \mathcal{B}^{\prime} \coloneqq \pi_{1} \cup \dots \cup \pi_{k-1} $. 
By definition, $ \mathcal{B}^{\prime} $ is MAT-free with MAT-partition $ (\pi_{1}, \dots, \pi_{k-1}) $. 
Also, the maximal exponents of $ \mathcal{B}^{\prime} $ are equal to $ k-1 $. 
Let $q:= \min \{k<i \le p\mid \pi_i \ne\varnothing \}$.
Since $ \pi_{q} \neq \varnothing $, we can take $ H \in  \pi_{q}$ and write $ \mathcal{B} \coloneqq \mathcal{B}^{\prime} \cup \{H\} $. 
It is a known fact\footnote{Here is the precise statement: ``Let  $ \mathcal{A} $ be an arrangement and let $ H \in \mathcal{A} $. If $ \mathcal{A}^{\prime} \coloneqq \mathcal{A}\setminus\{H\} $ is free with maximal exponent $ m $, then $ |\mathcal{A}^{\prime}| - |\mathcal{A}^{H}| \leq m $.'' 
We believe that this fact is well known among experts, but we give here a short proof for the sake of completeness. 
There exists a polynomial $ B $ such that $ \deg B = |\mathcal{A}^{\prime}| - |\mathcal{A}^{H}| $ and $ D(\mathcal{A}^{\prime})\alpha_{H} $ is contained in the ideal $(\alpha_{H}, B) \subseteq S$ \cite[Lemma 4.39 and Proposition 4.41]{OT92}. 
If $ \deg B > m $, then $ D(\mathcal{A}^{\prime})\alpha_{H} \subseteq (\alpha_{H}) $, hence $ D(\mathcal{A}^{\prime}) = D(\mathcal{A}) $. 
Therefore $ \mathcal{A} $ is free and  $\exp (\mathcal{A} )=\exp (\mathcal{A}^{\prime} )$. 
This implies $ |\mathcal{A}^{\prime}| = |\mathcal{A}| $, a contradiction. 
Thus $ |\mathcal{A}^{\prime}| - |\mathcal{A}^{H}| = \deg B \leq m $.}
 in the theory of free arrangements that $ |\mathcal{B}^{\prime}| - |\mathcal{B}^{H}| \leq k-1$.
However, \ref{definition MAT partition 3} implies $ |\mathcal{B}^{\prime}| - |\mathcal{B}^{H}| = q-1 >k-1$, a contradiction. 
This completes the proof.
\end{proof}

An \textbf{edge-labeled graph} is pair $ (G,\lambda) $ where $ G $ is a simple graph and  $\mbox{$\lambda \colon E_{G} \to \mathbb{Z}_{>0}$}$ is a map, called \textbf{(edge-)labeling}. 
Now we define a labeling of graphs which characterizes the MAT-freeness of graphic arrangements.
\begin{definition}[MAT-labelings]
\label{definition MAT-labeling}
Let $ (G,\lambda) $ be an  edge-labeled graph. 
Let $ \pi_{k} \coloneqq \lambda^{-1}(k) \subseteq E_{G}$ and $ E_{k} := \pi_{1} \sqcup \dots \sqcup \pi_{k} $ for every $ k \in \mathbb{Z}_{>0} $ and $ E_{0} \coloneqq \varnothing $. 
We say that $ \lambda $ is an \textbf{MAT-labeling} if the following conditions hold for every $ k \in \mathbb{Z}_{>0} $. 
\begin{enumerate}[(ML1)]
\item\label{definition MAT-labeling forest} $ \pi_{k} $ is a forest. 
\item\label{definition MAT-labeling closure} $ \cl(\pi_{k}) \cap E_{k-1} = \varnothing $. 
Here  $\cl(F)$ for $F \subseteq E_{G}$ denotes the \textbf{closure} of $F$ in the matroid sense. Namely,
an edge $e\in E_{G}$ is in $\cl(F)$ when the two endvertices of $e$ are connected by edges in $F$.
\item\label{definition MAT-labeling triangle} Every $ e \in \pi_{k} $ forms exactly $ k-1 $ triangles ($3$-cycles) with edges in $ E_{k-1} $. 
\end{enumerate}
\end{definition}

\begin{proposition}
\label{prop:MATF=MATL}
Let $G$ be a simple graph and $\K$ be a field. 
The graphic arrangement $\A_G$ in $\K^\ell$ is MAT-free if and only if $ G $ admits an MAT-labeling. 
\end{proposition}
\begin{proof}
This is a translation of Definition \ref{def:MAT-partition-free} into graphical terms with the use of Proposition \ref{no internal empty block}. 
\end{proof}

Thus characterizing the MAT-freeness of graphic arrangements amounts to characterizing the graphs having MAT-labelings. 
Here are first and simple facts on MAT-labelings.
Denote by $ K_{\ell} $ the complete graph on $ \ell $ vertices. 
\begin{proposition}\label{MAT-labeling on complete graph}
If $ \lambda $ is an MAT-labeling of $ K_{\ell} $, then $ |\pi_{k}| = \ell-k $ for all $ k \in [\ell-1] $. 
\end{proposition}
\begin{proof}
The graphic arrangement $ \mathcal{A}_{K_{\ell}} $ (in $\R^\ell$) is precisely the Weyl arrangement of type $A_{\ell-1}$ (also known as the \textbf{braid arrangement}) which has exponents $ \{0,1,2, \dots, \ell-1\}$.
Corollary \ref{cor:MAT} completes the proof.
\end{proof}

\begin{remark}
\label{rem:complete}
In fact, $ K_{\ell} $ always has an MAT-labeling $\lambda$. 
We can see this from Theorem \ref{thm:ideal-free} as $ K_{\ell} $ corresponds to a positive system (in particular, an ideal) of a root system of type $A$. 
More precisely,  $ \lambda \colon E_{K_{\ell}} \to \mathbb{Z}_{>0} $ is given by $ \lambda(\{v_{i}, v_{j}\}) = j-i$ for $1\le i < j \le \ell$ according to the height\footnote{The \textbf{height} of a positive root $\beta = \sum_{\alpha\in \Delta} c_\alpha \alpha \in \Phi^+$ is defined by $\sum_{\alpha\in \Delta} c_\alpha$.} 
 of positive roots. 
 Also, one can check directly that this labeling satisfies \ref{definition MAT-labeling forest}, \ref{definition MAT-labeling closure} and \ref{definition MAT-labeling triangle}.
We will see a different\footnote{Two edge-labeled graphs $(G_1,\lambda_1)$ and $(G_2, \lambda_2)$ are \textbf{isomorphic} if there exists a bijection $\sigma :V_{G_1} \to V_{G_2}$ such that $\{v_{i}, v_{j}\} \in E_{G_1}$ if and only if $\{\sigma(v_{i}), \sigma(v_{j})\} \in E_{G_2}$   with $ \lambda_1(\{v_{i}, v_{j}\}) =  \lambda_2(\{\sigma(v_{i}), \sigma(v_{j})\})$. If $G_1=G_2=G$, then we say that two labelings $\lambda_1$ and $ \lambda_2$ are the same (or isomorphic) if $(G,\lambda_1)$ and $(G, \lambda_2)$ are isomorphic.}  
(or nonisomorphic) labeling of $ K_{\ell} $ in Lemma \ref{MAT-L-PEO of complete graph}\ref{MAT-L-PEO of complete graph 2} (see also Figure \ref{fig:MAT-K4}). 

\end{remark}

It is important to know whether or not a restriction of an MAT-labeling is also an MAT-labeling. 
The proposition below states that it is enough to check the third condition.
\begin{proposition}\label{restriction to subset}
Let $ \lambda $ be an MAT-labeling of a simple graph $ G $ and $ F \subseteq E_{G} $. 
Then the restriction $ \lambda|_{F} $ is an MAT-labeling of the subgraph $ G_{F}= (V_{G}, F) $ if and only if~$ \lambda|_{F} $ satisfies \ref{definition MAT-labeling triangle}. 
\end{proposition}
\begin{proof}
Since $ (\lambda|_{F})^{-1}(k) = \pi_{k} \cap F \subseteq \pi_{k} $ and $ \pi_{k} $ is a forest, $ (\lambda|_{F})^{-1}(k) $ is also a forest. 
Moreover, $ \cl_{G_{F}}(\pi_{k} \cap F) \cap (E_{k-1} \cap F) \subseteq \cl_{G}(\pi_{k}) \cap E_{k-1} = \varnothing $. 
Thus  \ref{definition MAT-labeling forest} and \ref{definition MAT-labeling closure} are automatically satisfied. 
\end{proof}

\begin{lemma}\label{restriction to the union}
Let $ \lambda $ be an MAT-labeling of a simple graph $ G $ and $ F_{1}, F_{2} \subseteq E_{G} $. 
If $ \lambda|_{F_{1}} $ and $ \lambda|_{F_{2}} $ are MAT-labelings, then $ \lambda|_{F_{1} \cup F_{2}} $ is an MAT-labeling. 
\end{lemma}
\begin{proof}
For every subset $ F \subseteq E_{G} $ and $ k \in \mathbb{Z}_{>0} $, let $ \pi_{k}^{F} \coloneqq \lambda|_{F}^{-1}(k) = \pi_{k} \cap F $.
By Proposition  \ref{restriction to subset}, it suffices to prove that $ \lambda|_{F_{1} \cup F_{2}} $ satisfies \ref{definition MAT-labeling triangle}. 

Let $ e \in \pi_{k}^{F_{1}\cup F_{2}} = \pi_{k}^{F_{1}} \cup \pi_{k}^{F_{2}} $. 
Without loss of generality, we may assume $ e \in \pi_{k}^{F_{1}} $. 
Since $ \lambda|_{F_{1}} $ is an MAT-labeling, $ e $ forms at least $ k-1 $ triangles with edges in~$ E_{k-1} \cap (F_{1} \cup F_{2}) $. 
Moreover, since $ \lambda $ is an MAT-labeling, $ e $ forms at most $ k-1 $ triangles with edges in $ E_{k-1} \cap (F_{1} \cup F_{2}) $. 
Therefore $ e $ forms exactly $ k-1 $ triangles with edges in~$ E_{k-1} \cap (F_{1} \cup F_{2}) $. 
\end{proof}

